\phantomsection
\addcontentsline{toc}{subsection}{Inciso c)}
\subsection*{Inciso c)}

% Seccion con muchos nombres de clase largos, que no se parten con guion.
\sloppy

\textit{Integra con Composite para manejar grupos de cuentas (ej. ``Cuentas de una sucursal'')
y con Iterator para recorrer las transacciones de una cuenta. ¿Cómo afecta el Composite a la
ejecución de operaciones sobre grupos?}

\vspace{0.4cm}

\subsubsection*{Composite: el grupo es un elemento más}

\texttt{GrupoCuentas} implementa \texttt{ElementoContable}, de modo que una sucursal se visita
con la misma llamada que una cuenta suelta. El grupo no sabe qué operación se le aplica: solo
propaga la visita a sus hijos y combina los resultados parciales.

{\small
\begin{verbatim}
public final class GrupoCuentas implements ElementoContable,
        Iterable<Movimiento> {

    private final List<ElementoContable> hijos =
            new CopyOnWriteArrayList<>();

    @Override
    public <R> R accept(VisitanteOperacion<R> v) {
        R acumulado = v.neutro();
        for (ElementoContable hijo : hijos) {
            R parcial = hijo.accept(v);
            acumulado = acumulado == null
                    ? parcial : v.combinar(acumulado, parcial);
        }
        return acumulado;
    }
}
\end{verbatim}
}

El reparto de responsabilidades queda nítido: \textbf{el compuesto decide el recorrido} y
\textbf{el visitante decide la operación y cómo se agregan los resultados}. Con esto, el
cliente escribe \texttt{elemento.accept(v)} sin saber si detrás hay una cuenta o un árbol de
cuatro niveles, que es justamente el propósito de Composite.

\subsubsection*{El peaje que el Composite le cobra al Visitor}

Añadir el compuesto significa añadir un \emph{tipo de elemento}, y eso es exactamente lo que
Visitor hace caro. Había dos caminos:

\begin{enumerate}[leftmargin=1.2cm]
    \item \textbf{El clásico del catálogo GoF:} declarar \texttt{visitGrupo(GrupoCuentas)} en
          la interfaz del visitante. Da al visitante el control del recorrido (puede podar
          ramas, cambiar el orden o tratar el grupo como una unidad), pero obliga a
          implementar ese método en los $V$ visitantes existentes.
    \item \textbf{El adoptado aquí:} el recorrido vive en el compuesto y la interfaz del
          visitante solo gana dos métodos \texttt{default}, \texttt{combinar} y
          \texttt{neutro}, con comportamiento razonable por omisión. Los cuatro visitantes ya
          escritos siguieron compilando sin un solo cambio.
\end{enumerate}

{\small
\begin{verbatim}
public interface VisitanteOperacion<R> {
    // ... los cuatro metodos de visita del inciso (b)

    default R combinar(R acumulado, R siguiente) {
        return siguiente;
    }

    default R neutro() {
        return null;
    }
}
\end{verbatim}
}

Conviene decirlo sin adornos: la interfaz del visitante \emph{sí} se tocó. Los métodos
\texttt{default} de Java bajan el coste de $V$ modificaciones a una sola, pero no lo anulan, y
la contrapartida es que el visitante pierde el control del recorrido. Esta es la asimetría que
el inciso (d) cuantifica, observada aquí en un caso real.

\subsubsection*{La combinación tiene que ser asociativa}

Si los resultados parciales se agregan, la función que los agrega debe ser asociativa, porque
de otro modo el resultado dependería del orden del recorrido y el grupo no podría visitarse en
paralelo. Las tres operaciones se comportan de forma distinta a este respecto:

\begin{itemize}[leftmargin=1.2cm]
    \item \texttt{CalcularInteresesVisitor} forma el monoide $(0, +)$: neutro $0$ y suma
          asociativa. El interés de una sucursal es la suma del de sus cuentas.
    \item \texttt{AuditarTransaccionVisitor} consolida informes: suma los asientos revisados,
          concatena los hallazgos etiquetados con su cuenta y se queda con el peor nivel de
          riesgo, siendo BAJO $<$ MEDIO $<$ ALTO. También asociativa.
    \item \texttt{EjecutarTransaccionVisitor} \emph{no} tiene combinación natural: ejecutar no
          produce un valor que agregar, sino un efecto en cada cuenta. Se queda con el
          comportamiento por omisión y eso es una señal de diseño, no un descuido: las
          operaciones con efectos no se agregan, se aplican, y su problema es el alcance de la
          atomicidad.
\end{itemize}

\subsubsection*{Iterator: recorrer las transacciones}

\texttt{Cuenta} implementa \texttt{Iterable<Movimiento>}, así que el cliente recorre los
asientos sin conocer la estructura interna y sin que la cuenta ceda su lista:

{\small
\begin{verbatim}
public abstract class Cuenta implements ElementoContable,
        Iterable<Movimiento> {

    private final List<Movimiento> historial =
            new CopyOnWriteArrayList<>();

    @Override
    public Iterator<Movimiento> iterator() {
        return historial.iterator();   // instantanea, no bloquea
    }
}
\end{verbatim}
}

La elección de la estructura interna no es neutral frente a la restricción de concurrencia. Al
ser una lista \textit{copy-on-write}, el iterador trabaja sobre una instantánea: no lanza
\texttt{ConcurrentModificationException} aunque otro hilo registre asientos durante el
recorrido, y no bloquea a quien escribe. El precio es que el recorrido es
\textit{weakly consistent}: puede no ver los asientos añadidos después de empezar. Para una
auditoría es la garantía correcta, y es mejor que la alternativa de sostener el cerrojo de la
cuenta durante todo el recorrido.

Para el árbol, \texttt{IteradorMovimientosGrupo} aplana la jerarquía en profundidad con una
pila explícita, sin materializar una lista con todos los movimientos. Así
\texttt{GrupoCuentas} también es \texttt{Iterable<Movimiento>} y la uniformidad entre hoja y
compuesto se mantiene en los dos ejes: se visita igual y se recorre igual.

\subsubsection*{Cómo afecta el Composite a la ejecución de operaciones sobre grupos}

La pregunta del inciso tiene cinco respuestas, y solo la primera es gratis.

\textbf{1. El cliente se simplifica.} Una sola llamada sirve para una cuenta, una sucursal o el
banco entero. Las operaciones del inciso (b) funcionaron sobre el árbol sin modificarse.

\textbf{2. El alcance deja de ser constante.} \texttt{accept} sobre una hoja es una operación;
sobre un grupo son $N$. Una operación cara deja de ser inocente cuando el cliente ya no ve
cuántas cuentas hay debajo, y conviene prever poda o límites de profundidad.

\textbf{3. La atomicidad no se hereda, y esto es lo central.} El decorador transaccional toma
el cerrojo de \emph{cada} cuenta al visitarla, así que da atomicidad \textbf{por cuenta} y no
por grupo. En la prueba 5 de la ejecución se ve el efecto: un retiro de 4\,000,00 aplicado
sobre el árbol se ejecutó en cuatro cuentas y fue rechazado en \texttt{COR-001} por exceder su
línea de sobregiro. El grupo quedó \textbf{parcialmente aplicado}. Hay dos salidas, con costes
distintos:

\begin{itemize}[leftmargin=1.2cm]
    \item \textbf{Tomar todos los cerrojos del grupo}, siempre en orden total por
          identificador para no abrir la puerta al interbloqueo:
\end{itemize}

{\small
\begin{verbatim}
public static <R> R atomica(GrupoCuentas grupo,
                            VisitanteOperacion<R> v) {
    List<Cuenta> cuentas = hojas(grupo);  // ordenadas por id
    int tomados = 0;
    try {
        for (Cuenta c : cuentas) {
            c.cerrojo().lock();
            tomados++;
        }
        return grupo.accept(v);
    } finally {
        for (int i = tomados - 1; i >= 0; i--) {
            cuentas.get(i).cerrojo().unlock();
        }
    }
}
\end{verbatim}
}

\begin{itemize}[leftmargin=1.2cm]
    \item \textbf{Aceptar la atomicidad por cuenta y compensar}, con el Command del inciso (a)
          registrando cada aplicación para poder deshacerla si una hoja falla: el esquema de
          las \textit{sagas}. Escala mucho mejor y renuncia a la consistencia instantánea.
\end{itemize}

Y una precisión que conviene no pasar por alto: tomar todos los cerrojos da
\textbf{aislamiento}, no atomicidad. Nadie más opera en esas cuentas mientras dura, pero si la
operación falla a mitad de camino las aplicaciones anteriores siguen hechas. Para
todo\,o\,nada hace falta además una pasada de validación previa, o la compensación de la
segunda vía. El patrón no regala la consistencia; solo la hace expresable.

\textbf{4. El interbloqueo se vuelve una amenaza real.} Dos operaciones de grupo cuyas cuentas
se solapan pueden tomarse los cerrojos en órdenes opuestos. El orden total por identificador,
que en el inciso (b) hacía falta solo para las transferencias de dos cuentas, pasa a ser
obligatorio para cualquier operación de grupo.

\textbf{5. El paralelismo sale gratis, pero solo con visitantes puros.} Como los visitantes
devuelven su resultado en lugar de acumularlo y la combinación es asociativa, el recorrido se
reparte entre hilos y se reduce al final. En la prueba 6 el recorrido secuencial y el paralelo
dan el mismo total. Con un visitante que acumulara en un campo, el mismo recorrido sería una
condición de carrera: el Composite multiplica tanto el paralelismo disponible como el daño de
un visitante mal diseñado.

\subsubsection*{Evidencia de ejecución}

Salida literal de \texttt{DemoComposite}:

{\small
\begin{verbatim}
1. La estructura: hojas y compuestos tras la misma interfaz
-----------------------------------------------------------
  + Banco
    + Sucursal Centro
      - AHO-001     Ahorro      saldo=    4,700.00
      - COR-001     Corriente   saldo=     -205.00
      + Banca de empresa
        - EMP-001     Empresarial saldo=   84,997.50
        - EMP-002     Empresarial saldo=  240,000.00
    + Sucursal Este
      - INV-001     Inversion   saldo=   18,980.00

2. Uniformidad: el mismo visitante sobre una hoja y sobre el arbol
------------------------------------------------------------------
  hoja  -> AHO-001     Ahorro      riesgo=BAJO   asientos=  1  sin hallazgos
  grupo -> GRUPO riesgo=ALTO asientos=5
           COR-001: sobregiro activo por 205.00
           EMP-001: 1 operacion reportable
           EMP-002: 1 operacion reportable
           INV-001: penalizaciones acumuladas: 20.00

3. Reduccion: un solo visitante, cualquier nivel del arbol
----------------------------------------------------------
  Banca de empresa   interes anual = 3,899.97
  Sucursal Centro    interes anual = 3,923.47
  Sucursal Este      interes anual = 759.20
  Banco              interes anual = 4,682.67
  suma de las hojas por separado = 4,682.67

4. Iterator: transacciones de una cuenta y de todo el arbol
-----------------------------------------------------------
  movimientos de EMP-001:
    P1         PAGO_SERVICIOS    -15,000.00  reportable al regulador
  recorrido aplanado del arbol: 5 asientos, salidas por 76,700.00

5. Alcance de la atomicidad al operar sobre un grupo
----------------------------------------------------
  retiro de 4,000.00 sobre el arbol, cuenta por cuenta:
    AHO-001         4,700.00 ->       700.00  aplicada
    COR-001          -205.00 ->      -205.00  RECHAZADA
    EMP-001        84,997.50 ->    80,997.50  aplicada
    EMP-002       240,000.00 ->   236,000.00  aplicada
    INV-001        18,980.00 ->    14,900.00  aplicada
  comision de 10.00 con todos los cerrojos tomados:
    AHO-001           700.00 ->       690.00  aplicada
    COR-001          -205.00 ->      -220.00  aplicada
    EMP-001        80,997.50 ->    80,985.00  aplicada
    EMP-002       236,000.00 ->   235,987.50  aplicada
    INV-001        14,900.00 ->    14,889.80  aplicada

6. Recorrido en paralelo del grupo con reduccion al final
---------------------------------------------------------
  secuencial = 4,402.71
  en paralelo = 4,402.71
\end{verbatim}
}

La prueba 3 incluye una comprobación que vale la pena mirar: el interés del nodo raíz coincide
con la suma de los intereses de las hojas calculados por separado, lo que verifica que la
reducción del compuesto no pierde ni duplica hojas.

\subsubsection*{Dos limitaciones del arreglo}

\textbf{La visita deja de corresponderse una a una con los elementos.} Un visitante que
contara visitas tendría que decidir si los grupos cuentan, y la interfaz no le obliga a
plantearse la pregunta. Es la clase de ambigüedad que conviene documentar en el contrato del
visitante.

\textbf{El recorrido del compuesto es recursivo.} \texttt{accept} desciende por la pila de
llamadas, de modo que una jerarquía patológicamente profunda la desbordaría. El iterador, en
cambio, está escrito con una pila explícita y no tiene ese límite: una asimetría a tener
presente si las jerarquías crecen.

\fussy
